
Standard RLHF evaluation benchmarks measure preference agreement but cannot assess preference diversity due to structural limitations in dataset design. We investigate whether preference entropy---Shannon entropy H = -$\Sigma$ p\_i log(p\_i) applied to human preference distributions---can serve as a bidirectional alignment metric, measuring not only whether AI aligns to humans but whether humans preserve critical evaluation capacity when interacting with aligned AI.

Our validation on Anthropic-HH (n=100 prompts) reveals that pairwise comparison formats, optimized for annotation cost-efficiency, are structurally incompatible with per-prompt entropy measurement. Entropy computation succeeded technically (100\% success rate, all values within theoretical bounds [0, ln(2)] nats) but yielded constant H = ln(2) $\approx$ 0.693 nats across all prompts (variance = 0). Root cause analysis confirms that each pairwise example compares different response candidates (A1 vs B1, A2 vs B2, ...), creating tautological 50/50 distributions when aggregated, rather than actual preference distributions over fixed response sets.

We document a dataset structure trade-off: pairwise-unique formats enable efficient reward model training but fundamentally cannot measure diversity; multi-annotator-fixed formats enable entropy measurement at higher annotation cost. This explains why existing RLHF benchmarks (InstructGPT, Constitutional AI, WebGPT) structurally cannot retrospectively measure diversity even if desired.

Our contributions include: (1) introducing preference entropy as a proposed bidirectional alignment proxy (untested empirically due to dataset limitations), (2) documenting the cost-measurement trade-off in dataset structure (validated on Anthropic-HH with format-based inference for other datasets), (3) validating entropy computation feasibility when data structure supports it (100\% technical success), and (4) proposing alternative measurement proxies when multi-annotator data is unavailable (response diversity entropy, intra-annotator variance, entropy-regularized RLHF---all untested).

The hypothesis mechanism (preference entropy collapses faster than performance improves under extended RLHF training, particularly on subjective tasks) remains untested due to dataset incompatibility, not falsified. Future work requires either multi-annotator benchmark subsets (1K prompts $\times$ 20 annotators, ~20K labels) or alternative proxies to test whether RLHF inadvertently imposes preference monoculture on tasks where diversity is legitimate.

\textbf{Keywords:} RLHF evaluation, preference diversity, bidirectional alignment, Shannon entropy, dataset structure, annotation trade-offs

